\documentclass[10pt,a4paper]{amsart}
\usepackage[margin=19mm]{geometry}
\usepackage{amsmath,amssymb,mathtools}
\usepackage[scr=rsfs,cal=euler]{mathalfa}
\usepackage{xfrac}
\usepackage[unicode,hidelinks,bookmarks=false]{hyperref}
\newcommand{\ie}{i.e.\ }
\newcommand{\cf}{cf.\ }
\newcommand{\resp}{respectively\ }
\newcommand{\Subsection}{\subsection}
\providecommand{\nobreakdash}{\nobreak\hbox{-}}
\setlength{\emergencystretch}{2em}
\multlinegap0pt
\usepackage{bm}
\usepackage{bbm}
\usepackage{dsfont}
\usepackage{xspace}
\usepackage{cancel}
\usepackage{mathtools}
\usepackage{amsthm}
\makeatletter
\@ifundefined{theorem}{\newtheorem{theorem}{Theorem}}{}
\@ifundefined{lemma}{\newtheorem{lemma}[theorem]{Lemma}}{}
\@ifundefined{proposition}{\newtheorem{proposition}[theorem]{Proposition}}{}
\makeatother
\DeclareMathOperator{\en}{End}
\newcommand{\C}{\mathbb{C}}
\newcommand{\Qal}{\overline{\mathbb{Q}}}
\newcommand{\Q}{\mathbb{Q}}
\newcommand{\abs}[1]{\left| #1 \right|}
\newcommand{\cc}{\mathbf{c}}
\newcommand{\eps}{\varepsilon}
\newcommand{\m}{\mu}
\newcommand{\pg}[1]{\left\lbrace #1 \right\rbrace}
\newcommand{\pt}[1]{\left( #1 \right)}
\renewcommand{\Im}{\mathrm{Im}}
\renewcommand{\Re}{\mathrm{Re}}
\renewcommand{\a}{\alpha}
\renewcommand{\b}{\beta}
\renewcommand{\c}{\mathcal{C}}
\renewcommand{\l}{\lambda}
\renewcommand{\r}{\rho}
\title[Isogeny relations in products of families of elliptic curves]{Isogeny relations in products of families of elliptic curves}
\author{Luca Ferrigno}
\date{Source: arXiv:2409.01408v3 (CC BY 4.0)}
\begin{document}
\maketitle

\noindent\emph{Source and licence.} Isogeny relations in products of families of elliptic curves. Version arXiv:2409.01408v3 (CC BY 4.0), distributed under the Creative Commons Attribution 4.0 International licence, as recorded in the arXiv licence metadata for this exact version and shown on the abstract page https://arxiv.org/abs/2409.01408v3. Full rights notice: licenses/arxiv-2409.01408-CC-BY-4.0.txt.\par\smallskip

\noindent\textit{Editorial scope.} Counted unit: the Proposition whose source label is \texttt{prop:main\_estim\_isog} in the article (source0.tex lines 448--450, source label \texttt{prop:main\_estim\_isog}), delivered together with the article's own complete original proof of it (source0.tex lines 499--530, one \texttt{proof} environment of 32 lines). No mathematics is written, repaired or re-derived: statement and proof are the article's own text line for line. The proof is detached from the statement in the source: the article states the result at lines 448--450 and prints its proof at lines 499--530, and the proof environment's own header names the statement, so the association is the article's own. Required context retained uncounted: thm:main\_thm\_isog (lines 149--152); lemma:bound\_ll\_Teps\_isog (lines 471--473). Every definition, display and label of those lines is retained unchanged. Omitted from this package and not redistributed: 1--148, 153--447, 451--470, 474--498, 531--698. Nothing inside a retained excerpt is omitted or altered: every retained body is delivered exactly as the article's lines are written, so no omission receipt of any kind accompanies this package. Citations: the retained excerpts carry no citation command at all, so no bibliography entry of the article is transcribed and no reference is redistributed. Reference adaptation: none was needed and none was made; every reference inside the delivered unit resolves inside it, so no unresolved reference or citation is produced. Printed slips kept literal: no printed word, symbol, hypothesis or number of the counted statement or of its proof was altered. Layout and class: the article's own class is \texttt{amsart} with packages including inputenc, babel, mathtools, amsmath, amsthm, amssymb, amsfonts, enumitem, graphicx, hyperref, geometry, xcolor; the wrapper is the lane's pinned \texttt{amsart} editorial wrapper at 10pt on A4, which restates the theorem-like environments the retained text needs and adds the article's own macro definitions that the retained text needs (\texttt{\textbackslash C}, \texttt{\textbackslash Im}, \texttt{\textbackslash Q}, \texttt{\textbackslash Qal}, \texttt{\textbackslash Re}, \texttt{\textbackslash a}, \texttt{\textbackslash abs}, \texttt{\textbackslash b}, \texttt{\textbackslash c}, \texttt{\textbackslash cc}, \texttt{\textbackslash en}, \texttt{\textbackslash eps}, \texttt{\textbackslash l}, \texttt{\textbackslash m}, \texttt{\textbackslash pg}, \texttt{\textbackslash pt}, \texttt{\textbackslash r}), with extra wrapper packages (bm, bbm, dsfont, xspace, cancel, mathtools). The delivered numbering is the unit's own. Assets: no retained span contains a figure environment, a graphics-inclusion command, a PDF-page inclusion, a table or a TikZ picture; no image file is copied into this package, the package declares no asset dependency and no image file is redistributed with it. Licence: version arXiv:2409.01408v3 of this article is distributed under the Creative Commons Attribution 4.0 International licence, as recorded in the arXiv licence metadata for this exact version and shown on the abstract page https://arxiv.org/abs/2409.01408v3. A full rights notice is delivered at licenses/arxiv-2409.01408-CC-BY-4.0.txt. Characters: every retained excerpt is pure ASCII, so the delivered engine, XeLaTeX with its default Latin Modern text font, reports no missing character.
\begin{theorem}\label{thm:main_thm_isog}
Let $\c \subseteq E_{\l}^m \times E_{\m}^n$ be an irreducible asymmetric curve defined over $\Qal$ not contained in a fixed fiber, and define $P_i, Q_j$ as above.
Suppose moreover that $E_{\l}$ and $E_{\m}$ are not generically isogenous on $\c$ and that there are no generic non-trivial relations among $P_1, \ldots, P_m$ on $E_{\l}$ and among $Q_1, \ldots, Q_n$ on $E_{\m}$ with coefficients in $R_1$ and $R_2$, respectively. Then, there are at most finitely many $\cc \in \c (\C)$ such that there exist an isogeny $\phi: E_{\m(\cc)} \rightarrow E_{\l(\cc)}$ and $(a_1, \ldots, a_{m+n}) \in \en(E_{\l(\cc)})^{m+n}\setminus \pg{\mathbf{0}}$ with $$a_1 P_1(\cc)+ \ldots +a_m P_m(\cc)+ a_{m+1} \phi \pt{Q_1(\cc)}+ \ldots + a_{m+n} \phi \pt{Q_n(\cc)}=O.$$
\end{theorem}

\begin{lemma}\label{lemma:bound_ll_Teps_isog}
For every $\eps>0$, $\# \pi_2\pt{W^{\sim}(\Q,T)}\ll_{\eps} T^{\eps}$, for all $T\geq 1$.
\end{lemma}

\begin{proposition}\label{prop:main_estim_isog}
Under the hypotheses of Theorem \ref{thm:main_thm_isog}, for all $\eps>0$, we have $\# \mathcal{Z}(T)\ll_{\eps} T^{\eps}$, for all $T\geq 1$.
\end{proposition}

\begin{proof}[Proof of Proposition \ref{prop:main_estim_isog}]
If $(\tau_1,\widetilde{z}_1,\ldots,\widetilde{z}_{\ell}, z_{\ell+1}, \ldots, z_m, \tau_2, w_1, \ldots, w_n) \in \mathcal{Z}(T)$, then there are integers $ a_1, \ldots, a_{m+n},b_1, \ldots, b_{m+n}$, $A,B,C,D $ with absolute value bounded by $T$ and integers $\xi_1, \xi_2$ such that 
\[ \pt{\a_{\ell+1}+\b_{\ell+1} \r, \ldots,\a_{m+n}+ \b_{m+n} \r}\neq \mathbf{0} \qquad AD - BC \neq 0  \qquad C\tau_1 + D \neq 0 \]
\[ \pt{C\tau_1+D}\tau_2=A\tau_1+B \]
\[ \sum\limits_{i=1}^{\ell}(a_i+b_i \r)\widetilde{z}_i + \sum\limits_{i=\ell+1}^{m}(a_i+b_i \r)z_i+\pt{C\tau_1+D}\sum\limits_{j=1}^{n}(a_{m+j}+b_{m+j}\r)w_j=\xi_1+\xi_2 \tau_1\]
And since $ \lvert \tau_1 \rvert, \lvert \tau_2 \rvert, \lvert A \rvert, \lvert B \rvert, \lvert C \rvert, \lvert D \rvert, \lvert a_{1} \rvert, \ldots,  \lvert a_{m+n} \rvert, \abs{b_{1}}, \ldots, \abs{b_{m+n}}\leq T$ and $\widetilde{z}_p, z_{q} \in \mathcal{L}_{\tau_1}$, $w_r \in \mathcal{L}_{\tau_2}$ we have that
\begin{align*}
&\left| \sum\limits_{p=1}^{\ell}(a_p+b_p \r)\widetilde{z}_p + \sum\limits_{q=\ell+1}^{m} (a_q +b_q \r)z_q + (C\tau_1+D)\sum\limits_{r=1}^{n} (a_{m+r}+b_{m+r}\r) w_r \right|\\
 &\leq \sum\limits_{p=1}^{\ell} \pt{\lvert a_p \rvert+ \abs{b_p} \abs{\r}}  \lvert \widetilde{z}_p \rvert + \sum\limits_{q=\ell+1}^{m} \pt{\lvert a_q \rvert+ \abs{b_q} \abs{\r}}  \lvert z_q \rvert + \lvert C\tau_1+D \rvert \sum\limits_{r=1}^{n} \pt{\lvert a_{m+r} \rvert +\abs{b_{m+r}}\abs{\r}} \lvert w_r \rvert \\
&\ll T \cdot \max \pg{1, \lvert \tau_1 \rvert}+ (T\cdot \max\pg{1, \abs{\tau_1}}) \cdot T \cdot \max \pg{1, \lvert \tau_2 \rvert} \ll T^4
\end{align*}
Therefore, we have 
$$\abs{\xi_1+ \xi_2\tau_1}=\abs{\sum\limits_{p=1}^{\ell}(a_p+b_p \r)\widetilde{z}_p + \sum\limits_{q=\ell+1}^{m} (a_q +b_q \r)z_q + (C\tau_1+D)\sum\limits_{r=1}^{n} (a_{m+r}+b_{m+r}\r) w_r}\ll T^4$$
from which we deduce that $\abs{\xi_2}\ll T^5$, since $\Im(\tau_1)\geq \frac{1}{T}$ and
$$T^4 \gg \abs{\xi_1+ \xi_2\tau_1} \geq \abs{\Im(\xi_1+ \xi_2\tau_1)}=\abs{\xi_2 \Im(\tau_1)}\geq \frac{\abs{\xi_2}}{T}.$$
It follows that $$\abs{\xi_1}=\abs{\xi_1+\xi_2\tau_1-\xi_2\tau_1}\leq \abs{\xi_1+\xi_2\tau_1}+\abs{\xi_2}\cdot \abs{\tau_1} \ll T^4+ T^5 \cdot T \ll T^6.$$
Hence we get that 
\begin{align*}
& \left( a_1, \ldots, a_{m+n},b_1, \ldots, b_{m+n}, A,B,C,D, \xi_1, \xi_2,\right. \\
& \left. \Re(\tau_1), \Im(\tau_1), \Re(\widetilde{z}_1), \Im(\widetilde{z}_1), \ldots, \Re(\widetilde{z}_{\ell}), \Im(\widetilde{z}_{\ell}), \Re(z_{\ell+1}), \Im(z_{\ell+1}), \ldots, \Re(z_m), \Im(z_m), \right.\\  
& \left. \Re(\tau_2), \Im(\tau_2), \Re(w_1), \Im(w_1), \ldots, \Re(w_n), \Im(w_n)\right) \in W^{\sim} (\Q, \nu T^6)
\end{align*} 
for some positive constant $\nu$.  
Finally, consider the map 
$$\setlength{\arraycolsep}{0pt}
\renewcommand{\arraystretch}{1.2}
  \begin{array}{ r c l }
    \mathcal{Z}(T) & {} \longrightarrow {} & \pi_2\pt{W^{\sim} (\Q, \nu T^6)}\\
     (\tau_1, \ldots, w_n)      & {} \longmapsto {} & \pt{\Re(\tau_1), \Im(\tau_1), \ldots, \Re(w_n), \Im(w_n)}.
  \end{array}$$
Since this map is injective, the conclusion follows from Lemma~\ref{lemma:bound_ll_Teps_isog}.
\end{proof}

\end{document}
